\documentclass[10pt,a4paper]{amsart}
\usepackage[margin=19mm]{geometry}
\usepackage{amsmath,amssymb,mathtools}
\usepackage[scr=rsfs,cal=euler]{mathalfa}
\usepackage{xfrac}
\usepackage[unicode,hidelinks,bookmarks=false]{hyperref}
\newcommand{\ie}{i.e.\ }
\newcommand{\cf}{cf.\ }
\newcommand{\resp}{respectively\ }
\newcommand{\Subsection}{\subsection}
\providecommand{\nobreakdash}{\nobreak\hbox{-}}
\setlength{\emergencystretch}{2em}
\multlinegap0pt
\usepackage{bm}
\usepackage{float}
\usepackage{algorithm}\usepackage{algpseudocode}
\usepackage{amssymb}
\usepackage{mathtools}
\usepackage[shortlabels]{enumitem}
\usepackage{tikz}
\newtheorem{lem}{Lemma}
\newcommand{\blockdiag}{{\mathrm{block\text{-}diag}}}
\title{Non-degenerate Rigid Alignment in a Patch Framework}
\author{Dhruv Kohli, Gal Mishne, Alexander Cloninger}
\date{Source: arXiv:2303.11620v4 (CC BY 4.0)}
\begin{document}
\maketitle

\noindent\emph{Source and licence.} Non-degenerate Rigid Alignment in a Patch Framework. Version arXiv:2303.11620v4 (CC BY 4.0), distributed by arXiv under the Creative Commons Attribution 4.0 International licence, http://creativecommons.org/licenses/by/4.0/, as recorded in the arXiv licence metadata for this exact version and shown on the abstract page https://arxiv.org/abs/2303.11620v4. Full licence notice: \texttt{licenses/arxiv-2303.11620-CC-BY-4.0.txt}.\par\smallskip

\noindent\textit{Editorial scope.} Counted unit: the article's lem rev2:lem1 (source0.tex lines 2788--2791). The article's complete original proof of it is delivered with it (source0.tex lines 2792--2810, one \texttt{proof} environment of 19 lines). No mathematics is written, repaired or re-derived: the statement and the proof are the article's own lines, in the article's own notation, and no retained line is substituted or re-flowed.  No further source-local context is needed: every cross-reference of every retained line resolves inside the two delivered spans.  Nothing was omitted from any retained span, so the package carries no omission record.  No retained line contains a citation, so the wrapper carries no bibliography and no bibliography entry.  No retained line contains a figure environment, an image-inclusion command or drawing code, so no image is delivered and the package declares no asset dependency.  Editorial formatting only, in addition: an AMS \texttt{amsart} preamble that restates the article's own declarations and macros used by the retained lines, namely the environments \texttt{lem} on separate counters, together with the standard packages those lines need.  The retained environments print in this document as Lemma 1 in order of appearance, whereas the article prints the same items with the numbering of its own full document.  The complete native source of the article as distributed in the arXiv e-print for this exact version is bundled unchanged as \texttt{sources/138963/source0.tex}.
\begin{lem}
\label{rev2:lem1}
    Let $\mathbf{O}_1, \mathbf{O}_2 \in \mathbb{O}(d)^m$ and $\widetilde{\mathbf{O}}_1 = \pi(\mathbf{O}_1)$ and $\widetilde{\mathbf{O}}_2 = \pi(\mathbf{O}_2)$. Let $\widetilde{\mathbf{O}} = \pi(\mathbf{O})$ be on the minimal geodesic between $\widetilde{\mathbf{O}}_1$ and $\widetilde{\mathbf{O}}_2$. Suppose $\min_{\mathbf{Q} \in \mathbb{O}(d)}\left\|\mathbf{O}_j - \mathbf{S}^*\mathbf{Q}\right\|_F < \min\{2,\frac{2}{\pi}\zeta\delta(\mathbf{S}^*)\}$ for both $j =1, 2$. Then, $\min_{\mathbf{Q} \in \mathbb{O}(d)}\left\|\mathbf{O} - \mathbf{S}^*\mathbf{Q}\right\|_F < \min\{\pi, \zeta\delta(\mathbf{S}^*)\}$.
\end{lem}

\begin{proof}
Define $d(\cdot,\cdot)$ to be the geodesic distance on $\mathbb{O}(d)$ equipped with the Euclidean metric induced by the ambient space i.e. $\mathbb{R}^{d \times d}$. Then, from the properties of a quotient and a product manifold, for a general $\widetilde{\mathbf{S}} = \pi(\mathbf{S})$ it trivially follows that 
\begin{align}
 d_{\widetilde{g}}(\widetilde{\mathbf{S}},\widetilde{\mathbf{S}}^{*}) &= \min_{\mathbf{Q}\in \mathbb{O}(d)} d_{g}(\mathbf{S},\mathbf{S}^{*}\mathbf{Q}) \geq \min_{\mathbf{Q}\in \mathbb{O}(d)} \left\|\mathbf{S} -\mathbf{S}^{*}\mathbf{Q}\right\|_F. \label{rev2:eq7}
\end{align}
Moreover, if $\min_{\mathbf{Q} \in \mathbb{O}(d)}\left\|\mathbf{S} - \mathbf{S}^*\mathbf{Q}\right\|_F < 2$ (meaning $\mathbf{S}_i$ and $\mathbf{S}^*_i\mathbf{Q}^*$ where $\mathbf{Q}^*$ is the minimizer, lie in the same component of $\mathbb{O}(d)$ for all $i \in [1,m]$), then
\begin{align}
 d_{\widetilde{g}}(\widetilde{\mathbf{S}},\widetilde{\mathbf{S}}^{*})^2 &= \min_{\mathbf{Q}\in \mathbb{O}(d)} d_{g}(\mathbf{S},\mathbf{S}^{*}\mathbf{Q})^2 \leq d_{g}(\mathbf{S},\mathbf{S}^{*}\mathbf{Q}^*)^2 = \sum_{i=1}^{m} d(\mathbf{S}_i,\mathbf{S}^{*}_i\mathbf{Q}^*)^2\\
 &\leq  \frac{\pi^2}{4} \sum_{i=1}^{m}\left\|\mathbf{S}_i -\mathbf{S}^{*}_i\mathbf{Q}^*\right\|_F^2 = \frac{\pi^2}{4} \left\|\mathbf{S} -\mathbf{S}^{*}\mathbf{Q}^*\right\|_F^2.\\
 &= \min_{\mathbf{Q}\in \mathbb{O}(d)}\frac{\pi^2}{4} \left\|\mathbf{S} -\mathbf{S}^{*}\mathbf{Q}\right\|_F^2.
\end{align}
Here, the second inequality follows from the following fact: if $\mathbf{Q}_1, \mathbf{Q}_2 \in \mathbb{O}(d)$ and are in the same component, then $d(\mathbf{Q}_1,\mathbf{Q}_2)^2 \leq \frac{\pi^2}{4} \left\|\mathbf{Q}_1-\mathbf{Q}_2\right\|_F^2$. This is because, since $\mathbf{Q}_1^T\mathbf{Q}_2 \in \mathbb{SO}(d)$ and therefore we can write $\mathbf{Q}_1^T\mathbf{Q}_2 = \mathbf{V}\mathbf{R}\mathbf{V}^T$ where $\mathbf{V} \in \mathbb{O}(d)$,
$\mathbf{R} = \blockdiag(R(\theta_1),\ldots,R(\theta_k),\mathbf{I}_r)$
and $R(\theta_j) = \begin{bmatrix}\cos(\theta_j) & -\sin(\theta_j)\\ \sin(\theta_j) & \cos(\theta_j)\end{bmatrix}$. Here $2k + r = d$ and $\theta_j \in (-\pi,\pi]$. Also, $\bm{\Omega} = \text{Log}(\mathbf{Q}_1^T\mathbf{Q}_2)$ has nonzero eigenvalues $\{\pm \iota \theta_j\}_{j=1}^{k}$. Putting all together,
\begin{align}
 d(\mathbf{Q}_1,\mathbf{Q}_2)^2 = \left\|\bm{\Omega}\right\|_F^2 = 2\sum_{j=1}^{k}\theta_j^2 \leq \frac{\pi^2}{4} \sum_{j=1}^{k} 4(1-\cos(\theta_j)) = \frac{\pi^2}{4} \left\|\mathbf{I}_d - \mathbf{Q}_1^T\mathbf{Q}_2\right\|_F^2 = \frac{\pi^2}{4} \left\|\mathbf{Q}_1-\mathbf{Q}_2\right\|_F^2.
\end{align}
It follows that $d_{\widetilde{g}}(\widetilde{\mathbf{O}}_j,\widetilde{\mathbf{S}}^{*}) < \min\{\pi, \zeta \delta(\mathbf{S}^*)\}$ for both $j=1,2$. Then, since a geodesic ball of radius less then $\pi$ is geodesically convex in $\mathbb{O}(d)$, the same result holds for geodesic balls on $\mathbb{O}(d)^m$. Consequently, $d_{\widetilde{g}}(\widetilde{\mathbf{O}},\widetilde{\mathbf{S}}^{*}) < \min\{\pi, \zeta \delta(\mathbf{S}^*)\}$. From Eq.~(\ref{rev2:eq7}), it follows that $\min_{\mathbf{Q}\in \mathbb{O}(d)} \left\|\mathbf{O} -\mathbf{S}^{*}\mathbf{Q}\right\|_F < \min\{\pi,\zeta \delta(\mathbf{S}^*)\}$.
\end{proof}

\end{document}
